\documentclass[11pt]{article}
\usepackage[margin=1in]{geometry}
\usepackage{amsmath,amssymb,amsthm}
\usepackage[T1]{fontenc}
\usepackage{lmodern}
\usepackage{microtype}
\usepackage{xurl}
\usepackage[colorlinks=true,allcolors=blue]{hyperref}
\hypersetup{pdfauthor={Kaiwen Jin},pdftitle={AIM 450 and a uniform Hilbert transform estimate for Lipschitz directions}}
\newtheorem{theorem}{Theorem}
\newtheorem{corollary}[theorem]{Corollary}
\newcommand{\R}{\mathbb R}
\newcommand{\Schwartz}{\mathcal S}
\title{AIM 450 and a uniform Hilbert transform estimate\\
for Lipschitz directions}
\author{Kaiwen Jin\thanks{Contact: \href{mailto:yui.ui.siki@gmail.com}{yui.ui.siki@gmail.com}.}}
\date{October 12, 2026}
\begin{document}
\maketitle
\begin{abstract}
Stein's question asks for a uniform weak-$L^2$ bound for the short Hilbert
transform along a Lipschitz unit vector field in the plane, allowing dependence
on both coordinates. OpenAI's September 25, 2026 preprint states a stronger
uniform $L^2$ estimate for every inner truncation. We record its exact
implication for AIM problem 450: the symmetric principal value exists for
Schwartz inputs, and the requested level-set bound holds with absolute
constants. The passage uses cancellation at zero and Fatou's lemma. The main
uniform estimate and its proof are due to OpenAI.
\end{abstract}

The Hilbert transform along a varying direction integrates a function on a
different straight line at each point. Stein's question asks whether a
Lipschitz bound on the direction controls this singular integral uniformly
on sufficiently short segments. Lacey and Li formulate the weak-$(2,2)$
conjecture in the final preprint version of
\cite[Conjecture~1.5, equation~(1.6)]{LaceyLi}; the normalized form appears
as AIM problem 450 \cite{AIM}.

OpenAI's preprint \cite[Theorem~1.1]{OpenAI} states the uniform estimate
needed for this question, together with its principal-value and weak-type
consequences. This note compares that statement with the exact catalogue
target and spells out the passage to those consequences. The uniform
estimate and its proof are due to OpenAI.

For a unit vector field $v\colon\R^2\to S^1$ and $0<\varepsilon<a$,
write
\begin{equation}\label{eq:truncated}
 H^{\varepsilon}_{v,a}f(x)
 =\int_{\varepsilon<|t|<a}f(x-tv(x))\,\frac{dt}{t}.
\end{equation}
Here $S^1=\{u\in\R^2:|u|=1\}$, and the direction $v(x)$ is
held fixed while $t$ varies. The integral follows the straight line through
$x$ in direction $v(x)$. We use Lebesgue measure on $\R^2$ and allow
complex-valued Schwartz functions $f\in\Schwartz(\R^2)$.

The following is the estimate stated in Theorem~1.1 of OpenAI's preprint
\cite{OpenAI}, with the operator notation in \eqref{eq:truncated}.

\begin{theorem}[OpenAI]\label{thm:openai}
There are absolute constants $a_*\in(0,1/2)$ and $C_*<\infty$
such that, for every $v\colon\R^2\to S^1$ satisfying
\[
 |v(x)-v(y)|\le |x-y|\qquad(x,y\in\R^2)
\]
and every $f\in\Schwartz(\R^2)$,
\begin{equation}\label{eq:uniform}
 \sup_{0<\varepsilon<a_*}
 \bigl\|H^{\varepsilon}_{v,a_*}f\bigr\|_{L^2(\R^2)}
 \le C_*\|f\|_{L^2(\R^2)}.
\end{equation}
\end{theorem}

The supremum in \eqref{eq:uniform} is outside the norm. The estimate
applies to fields depending on both spatial coordinates and to unrestricted
Schwartz inputs. Its full proof is in the cited preprint; the argument below
uses precisely this stated estimate.

\begin{corollary}[The bound in AIM 450]\label{cor:aim}
There are absolute constants $K\ge1$ and $C<\infty$ such that, for every
$v\colon\R^2\to S^1$ with $|v(x)-v(y)|\le|x-y|$, every
$f\in\Schwartz(\R^2)$, and every $\lambda>0$, the symmetric principal value
\begin{equation}\label{eq:pv}
 H_vf(x)=\lim_{\varepsilon\downarrow0}
 \int_{\varepsilon<|t|<1/K}f(x-tv(x))\,\frac{dt}{t}
\end{equation}
exists at every $x\in\R^2$ and satisfies
\begin{equation}\label{eq:weak}
 \lambda^2\bigl|\{x\in\R^2:|H_vf(x)|>\lambda\}\bigr|
 \le C\|f\|_{L^2(\R^2)}^2.
\end{equation}
One may take $K=1/a_*$ and $C=C_*^2$, with the constants from
Theorem~\ref{thm:openai}.
\end{corollary}

\begin{proof}
Choose $a_*$ and $C_*$ from Theorem~\ref{thm:openai} and put
$K=1/a_*>2$. For fixed $x$, changing variables in the negative half of the
symmetric integral gives
\begin{equation}\label{eq:symmetric}
 H^{\varepsilon}_{v,a_*}f(x)
 =\int_{\varepsilon}^{a_*}
 \frac{f(x-tv(x))-f(x+tv(x))}{t}\,dt.
\end{equation}
Set $M_x=\sup_{|z-x|\le a_*}|\nabla f(z)|$, where the gradient norm is
the Euclidean norm of its complex components. Smoothness of $f$ gives
$M_x<\infty$. Since $|v(x)|=1$, the fundamental theorem of calculus along
the segment $x+s v(x)$ yields
\[
 |f(x-tv(x))-f(x+tv(x))|
 \le \int_{-t}^{t}|\nabla f(x+s v(x))|\,ds
 \le 2t M_x\qquad(0<t<a_*).
\]
Thus the integrand in \eqref{eq:symmetric} is bounded by $2M_x$ on
$(0,a_*)$. Dominated convergence proves the existence of the limit in
\eqref{eq:pv} at each $x$.

For fixed $\varepsilon>0$, continuity of $v$ and $f$ makes
$H^{\varepsilon}_{v,a_*}f$ continuous, by dominated convergence in $t$ on
the truncated interval. Hence the principal value is measurable.
Take $\varepsilon_n=a_* /(n+2)$ for $n=0,1,\ldots$. Fatou's lemma and
\eqref{eq:uniform} give
\[
 \int_{\R^2}|H_vf(x)|^2\,dx
 \le \liminf_{n\to\infty}
       \int_{\R^2}|H^{\varepsilon_n}_{v,a_*}f(x)|^2\,dx
 \le C_*^2\|f\|_{L^2(\R^2)}^2.
\]
Consequently, for each $\lambda>0$,
\[
 \lambda^2|\{x:|H_vf(x)|>\lambda\}|
 \le \int_{\R^2}|H_vf(x)|^2\,dx
 \le C_*^2\|f\|_{L^2(\R^2)}^2,
\]
which is \eqref{eq:weak}. Both constants were chosen before $v$, $f$,
and $\lambda$.
\end{proof}

The cutoff $1/K=a_*$ is fixed and absolute. Corollary~\ref{cor:aim}
retains the two-coordinate direction fields and every quantifier in the
pinned AIM statement \cite{AIM}. No estimate for a supremum inside the
$L^2$ norm, or for arbitrary outer lengths, is used here.

\paragraph{AI assistance.}
OpenAI's preprint \cite{OpenAI} supplies the uniform estimate and its proof.
This note was prepared with assistance from OpenAI's ChatGPT/Codex for
source comparison, drafting, and proof checking.

\begin{thebibliography}{9}
\bibitem{AIM}
AIM problem catalogue, \emph{450. Stein's Hilbert-transform bound along Lipschitz
directions}, AIM problem catalogue, revision
\texttt{c9929805cece044705abaaebe1cebddac1ba11a6}.
\url{https://github.com/MColbrook/AIM/blob/c9929805cece044705abaaebe1cebddac1ba11a6/problems/450-stein-lipschitz-direction-hilbert-bound.md}.

\bibitem{LaceyLi}
M. Lacey and X. Li, \emph{On a conjecture of E. M. Stein on the Hilbert
transform on vector fields}, Memoirs of the American Mathematical Society
\textbf{205} (2010), no.~965.
\url{https://doi.org/10.1090/S0065-9266-10-00572-7}.
Final preprint: arXiv:0704.0808v3, January~11, 2008; Conjecture~1.5 and
equation~(1.6). \url{https://arxiv.org/abs/0704.0808v3}.

\bibitem{OpenAI}
OpenAI, \emph{A uniform Hilbert transform estimate for Lipschitz directions},
preprint, September~25, 2026; Theorem~1.1. Revision
\texttt{adc7f1241b42e322a6451854ab7e4b4c146bf78a}.
\url{https://github.com/openai/math/blob/adc7f1241b42e322a6451854ab7e4b4c146bf78a/preprints/A-uniform-Hilbert-transform-estimate-for-Lipschitz-directions-September-25-2026/main.pdf}.
\end{thebibliography}
\end{document}
